\section*{Chapter \ref{ch:b1:observing-the-sky} --- Telescopes, Detectors and Photometry}

\begin{solution}{exo:b1:observing-the-sky:1}
A difference of 5 magnitudes is a factor of 100: the star of magnitude 1.0
sends 100 times more \omterm{def:b1:observing-the-sky:flux}{flux}. Rigel and Polaris: $10^{0.4\times(2.02 - 0.13)} =
10^{0.756} = 5.70$; Rigel, with the smaller magnitude, is brighter.
\end{solution}

\begin{solution}{exo:b1:observing-the-sky:2}
$d = 1/0.19495 = \qty{5.13}{\parsec}$, that is $5.1295\times206265 =
\qty{1.058e6}{\au}$ and $5.1295\times3.2616 = \qty{16.7}{\lightyear}$.
\end{solution}

\begin{solution}{exo:b1:observing-the-sky:3}
$1.22\times\num{550e-9}/2.4 = \qty{2.80e-7}{rad} = \ang{;;0.058}$, about
seven times finer than a \omterm{def:b1:observing-the-sky:seeing}{seeing} of \ang{;;0.40}: in space the telescope
reaches its diffraction limit. For \qty{6.5}{m} at \qty{2.0}{\micro\metre}:
$1.22\times\num{2.0e-6}/6.5 = \qty{3.75e-7}{rad} = \ang{;;0.077}$.
\end{solution}

\begin{solution}{exo:b1:observing-the-sky:4}
$B - V$ rises from Rigel ($-0.03$) to Altair (0.22), Polaris (0.60) and
Betelgeuse (1.85): the hottest is Rigel, then Altair, Polaris and
Betelgeuse. Distance dims both bands by the same factor, so $B - V$ is
unchanged by it (as long as no dust reddens the light on the way).
\end{solution}

\begin{solution}{exo:b1:observing-the-sky:5}
$X_1 = \sec\ang{30} = 1.1547$, $X_2 = \sec\ang{60} = 2.0000$.
$k = (6.590 - 6.412)/(2.0000 - 1.1547) = 0.178/0.8453 = \num{0.211}$ per
\omterm{def:b1:observing-the-sky:extinction}{airmass}, and $m_0 = 6.412 - 0.2106\times1.1547 = 6.169$.
\end{solution}

\begin{solution}{exo:b1:observing-the-sky:6}
$\mathrm{SNR} = \sqrt{400} = 20$. The SNR grows as $\sqrt t$: for 100,
$t = \qty{10}{s}\times(100/20)^2 = \qty{250}{s}$. An error of \num{0.01}
magnitude needs $\mathrm{SNR} = 1.086/0.01 = 108.6$, so
$t = 10\times(108.6/20)^2 = \qty{295}{s}$.
\end{solution}

\begin{solution}{exo:b1:observing-the-sky:7}
The distance is $d = 1/0.00378 = \qty{264.6}{\parsec}$, and the \omterm{def:b1:observing-the-sky:absolute}{absolute
magnitude} $M = 0.13 - 5\log_{10}26.46 = 0.13 - 7.11 = -6.98$. $L/L_\odot = 10^{0.4\times(4.83 + 6.98)} =
10^{4.724} = \num{5.3e4}$ in V.
\end{solution}

\begin{solution}{exo:b1:observing-the-sky:8}
$\pi\times1.2^2 = \qty{4.52}{m^2}$, and $25/4.52 = 5.5$: the Webb telescope
collects the same number of photons in an exposure 5.5 times shorter. A full
disc \qty{6.5}{m} across would have $\pi\times3.25^2 = \qty{33.2}{m^2}$; the
real mirror has $25/33.2 = \qty{75}{\percent}$ of it, because 18 hexagons
leave the edge of that disc uncovered and the gaps between segments and the
central opening collect nothing.
\end{solution}

\begin{solution}{exo:b1:observing-the-sky:9}
(a) $\qty{10}{\parsec} = \qty{3.0857e17}{m}$, and $F = \num{3.828e26}/(4\pi
\times(\num{3.0857e17})^2) = \num{3.828e26}/\num{1.1965e36} =
\qty{3.20e-10}{W.m^{-2}}$, $1361/\num{3.20e-10} = \num{4.3e12}$ times less than at
the Earth.
(b) $5\log_{10}(d/\qty{10}{\parsec}) = -26.74 - 4.83 = -31.57$, so
$d = \qty{10}{\parsec}\times10^{-6.314} = \qty{4.853e-6}{\parsec}$, that is
$\num{4.853e-6}\times206265 = \qty{1.001}{\au}$.
\end{solution}

\begin{solution}{exo:b1:observing-the-sky:10}
Above the atmosphere the \omterm{def:b1:observing-the-sky:flux}{flux} is $F_0$; below, $F = F_0e^{-\tau_0X}$. Then
$m - m_0 = -2.5\log_{10}(F/F_0) = 2.5\log_{10}e\;\tau_0X$, which is $kX$ with
$k = 2.5\log_{10}e\;\tau_0 = 1.0857\,\tau_0$. For $k = 0.160$:
$\tau_0 = 0.160/1.0857 = 0.147$. At the \omterm{def:b1:spherical-astronomy:celestial-sphere}{zenith} $e^{-0.147} = 0.863$ of the
light arrives, \qty{86}{\percent}; at $X = 2.31$, $e^{-0.147\times2.31} =
e^{-0.340} = 0.711$, \qty{71}{\percent}. At $z = \ang{85}$, $\sec z =
11.5$; but the plane-parallel model makes the path infinite at the horizon,
whereas the path through a curved atmosphere stays finite, and refraction
bends the ray as well: near the horizon the formula overestimates the
\omterm{def:b1:observing-the-sky:extinction}{airmass} and its errors grow fast.
\end{solution}

\begin{solution}{exo:b1:observing-the-sky:11}
(a) $\mathrm{SNR} = Rt/\sqrt{(R + B)t}$, so $t = \mathrm{SNR}^2(R + B)/R^2
= 100\times2050/2500 = \qty{82}{s}$. (b) Halving the \omterm{def:b1:observing-the-sky:seeing}{seeing} divides the
area of the image, hence $B$, by 4: $B = 500$ and $t = 100\times550/2500 =
\qty{22}{s}$. (c) For $B \gg R$, $t \approx \mathrm{SNR}^2B/R^2$. The source
rate $R$ grows as $D^2$; the sky rate $B$ grows as the collecting area
$D^2$ times the sky area of the image, $\theta^2$. Hence $t \propto
\theta^2D^2/D^4 = \theta^2/D^2$: halving the \omterm{def:b1:observing-the-sky:seeing}{seeing} is worth as much as
doubling the diameter.
\end{solution}

\begin{solution}{exo:b1:observing-the-sky:12}
(a) $N = 1200\times100 = \num{120000}$ lines; in the first order the
resolving power is \num{1.2e5}, and $\Delta\lambda = 550/(\num{1.2e5}) =
\qty{4.6e-3}{nm}$. (b) $v = c\,\Delta\lambda/\lambda = c/(\num{1.2e5}) =
\qty{2.5}{km.s^{-1}}$. (c) $\Delta\lambda = 550\times29.78/\num{299792} =
\qty{0.055}{nm}$, about 12 resolution elements. The observer must compute,
for the date and time of the observation, the component of the Earth's
orbital velocity toward the star, and subtract it, to give the star's
velocity relative to the Sun.
\end{solution}

\begin{solution}{pb:b1:observing-the-sky:1}
\textbf{1.} $d = 1/0.37921 = \qty{2.637}{\parsec}$.

\textbf{2.} $2.637\times206265 = \qty{5.44e5}{\au}$;
$2.637\times\num{3.0857e16} = \qty{8.14e16}{m}$.

\textbf{3.} $2.637\times3.2616 = \qty{8.60}{\lightyear}$: the light left
Sirius 8.6 years ago.

\textbf{4.} The direction swings by $2p = \ang{;;0.758}$, about twice the
\omterm{def:b1:observing-the-sky:seeing}{seeing}. The swing is measured from the position of the centre of the image,
not from its edge: averaged over many photons, the centre is found to a
small fraction of the image's width, the more precisely the more photons
are counted.

\textbf{5.} $\Delta d/d = \Delta p/p = 0.001/0.37921 = \qty{0.26}{\percent}$,
that is $8.60\times0.0026 = \qty{0.023}{\lightyear}$.

\textbf{6.} $10^{0.4\times(0.03 + 1.46)} = 10^{0.596} = 3.94$.

\textbf{7.} $10^{0.4\times(26.74 - 1.46)} = 10^{10.112} = \num{1.29e10}$.

\textbf{8.} $10^{0.4\times(8.44 + 1.46)} = 10^{3.96} = \num{9120}$.

\textbf{9.} $z = \ang{90} - \ang{82.09} = \ang{7.91}$, $X = \sec\ang{7.91} =
1.0096$; extinction $0.15\times1.0096 = 0.151$ magnitude; measured
$m = -1.46 + 0.151 = -1.31$.

\textbf{10.} $z = \ang{60}$, $X = 2.000$; extinction 0.30 magnitude,
$m = -1.16$. The fraction transmitted is $10^{-0.4\times0.30} = 0.759$:
\qty{24}{\percent} of the light is lost.

\textbf{11.} At rising $z = \ang{90}$ and $\sec z$ is infinite, while the
real path through the curved atmosphere is long but finite; refraction also
bends the ray. Near the horizon the extinction is large and changes fast,
and the plane-parallel correction is wrong.

\textbf{12.} $\num{3.0e4}\times10^{0.4\times(10.0 - 8.44)} = \num{3.0e4}
\times10^{0.624} = \num{3.0e4}\times4.21 = \num{1.26e5}$ photoelectrons per
second.

\textbf{13.} An SNR of 1000 needs $N = 1000^2 = \num{1.0e6}$
photoelectrons: $t = \num{1.0e6}/\num{1.26e5} = \qty{7.9}{s}$. The magnitude
error is $1.086/1000 = \num{0.0011}$.

\textbf{14.} $\num{3.0e4}\times10^{0.4\times(10.0 + 1.46)} = \num{3.0e4}
\times10^{4.584} = \num{1.15e9}$ photoelectrons per second.

\textbf{15.} $0.20\times\num{1.15e9} = \num{2.30e8}$ per second in the
brightest pixel, which fills in $\num{1.0e5}/\num{2.30e8} =
\qty{4.3e-4}{s}$, less than half a millisecond.

\textbf{16.} The wings of Sirius give $\num{1.15e9}\times10^{-4} =
\num{1.15e5}$ per second; Sirius~B gives $0.20\times\num{1.26e5} =
\num{2.52e4}$ per second. The glare of Sirius is 4.6 times brighter than the
peak of Sirius~B.

\textbf{17.} The pixel receives $\num{1.15e5} + \num{2.52e4} = \num{1.40e5}$
per second and fills in $\num{1.0e5}/\num{1.40e5} = \qty{0.71}{s}$. In
\qty{0.5}{s}: signal $\num{2.52e4}\times0.5 = \num{1.26e4}$, total
$\num{7.01e4}$, noise $\sqrt{\num{7.01e4}} = 265$, so $\mathrm{SNR} =
\num{1.26e4}/265 = 48$. Sirius~B is detectable as far as photon noise goes;
what makes it hard is that the glare must be known and subtracted to much
better than a fifth of its value, while the core of Sirius saturates in
under a millisecond.

\textbf{18.} $m - M = 5\log_{10}(2.637/10) = 5\times(-0.579) = -2.89$.

\textbf{19.} $M = -1.46 + 2.89 = 1.43$.

\textbf{20.} $d = 1/0.13023 = \qty{7.679}{\parsec}$;
$m - M = 5\log_{10}0.7679 = -0.57$, so $M = 0.03 + 0.57 = 0.60$.

\textbf{21.} Vega, with the smaller \omterm{def:b1:observing-the-sky:absolute}{absolute magnitude}, by the factor
$10^{0.4\times(1.43 - 0.60)} = 10^{0.332} = 2.15$. Sirius looks brighter
only because it is nearer.

\textbf{22.} $m = M_\odot + (m - M) = 4.83 - 2.89 = 1.94$: the Sun would be
a fairly bright naked-eye star.

\textbf{23.} $5\log_{10}(d/\qty{10}{\parsec}) = 6.0 - 1.43 = 4.57$, so
$d = 10\times10^{0.914} = \qty{82.0}{\parsec}$, about \qty{267}{\lightyear}.

\textbf{24.} $L/L_\odot = 10^{0.4\times(4.83 - 1.43)} = 10^{1.36} = 22.9$:
in the V band, Sirius outshines \textbf{about 23 suns}.
\end{solution}
